\section{商集与轨道完备性}\label{app:orbit-completeness}

本附录给出定理~\ref{thm:smallest-counterexample} 和定理~
\ref{thm:low-period-frontier} 所用的有限覆盖论证。

\subsection{切换坐标}

对 \(G_n=C_n(1,2)\)，\(n\) 个三角形圈和步长一 Hamilton 圈构成
\((n+1)\) 维圈空间的一组基。因此每个切换类由下式唯一表示：

\begin{equation}
\label{eq:cycle-coordinate-space}
(\tau_{0},\ldots,\tau_{n-1},\alpha) \in \{+1,-1\}^{n+1}.
\end{equation}

当 \(n\) 为偶数时，全局边符号取反会取反每个三角形通量并保持 \(\alpha\)；
它还把 \(A\) 换成 \(-A\)，故谱半径不变。相邻乘积
\(Q_i=\tau_i\tau_{i+1}\) 在该取反下不变。合法 \(Q\) 字的乘积为 \(+1\)，
且每个字恰有两个提升 \(\tau\) 和 \(-\tau\)。所以一对 \((Q,\alpha)\) 恰好
代表这一对谱等价切换类。

二面体群按图自同构作用。若规范 \(Q\) 轨道大小为 \(s\)，保留两个
\(\alpha\) 值就代表 \(4s\) 个切换类：\(s\) 个几何字、两个 \(\tau\) 提升
和两种和乐。对所有合法 \(Q\) 轨道求和得

\begin{equation}
\label{eq:orbit-weight-identity}
\sum_{\mathrm{orbits}} 4s = 4\cdot 2^{n-1}=2^{n+1},
\end{equation}

恰等于切换类总数。因此该商集不遗漏任何类。

\subsection{合法通量字的 Burnside 计数}

设循环的 \(p\) 个位置上的置换 \(g\) 具有圈长
\(\ell_1,\ldots,\ell_c\)。被 \(g\) 固定的字在每个圈上为常数，所以由符号
\(\varepsilon_1,\ldots,\varepsilon_c\) 指定。其总乘积为

\begin{equation}
\label{eq:permutation-cycle-sign}
\prod_{j=1}^c \varepsilon_j^{\ell_j}.
\end{equation}

若至少一个 \(\ell_j\) 为奇数，则 \(2^c\) 种赋值中恰有一半使式~
\eqref{eq:permutation-cycle-sign} 为正，所以 \(g\) 固定 \(2^{c-1}\) 个合法字。
若所有 \(\ell_j\) 均为偶数，则该式恒正，\(g\) 固定全部 \(2^c\) 种赋值。
Burnside 引理于是给出

\begin{equation}
\label{eq:burnside-orbit-count}
N_p=\frac{1}{2p}\sum_{g\in D_p}\begin{cases}2^{c(g)-1}, & \text{若 }g\text{ 有奇圈}, \\
2^{c(g)}, & \text{若 }g\text{ 的所有圈均为偶圈}
\end{cases}
\text{。}
\end{equation}

计算 \(p\) 个旋转和 \(p\) 个反射的圈分解，得到：

\begin{table}[tbp]
\centering
\caption{直至周期十六的合法通量字和二面体轨道数。}
\label{tab:low-period-orbit-counts}
\begingroup
\renewcommand{\arraystretch}{1.12}
\begin{tabular}{lll}
\toprule
\(p\) & 合法字数 & 二面体轨道数 \\
\midrule
1 & 1 & 1 \\
2 & 2 & 2 \\
3 & 4 & 2 \\
4 & 8 & 4 \\
5 & 16 & 4 \\
6 & 32 & 8 \\
7 & 64 & 9 \\
8 & 128 & 18 \\
9 & 256 & 23 \\
10 & 512 & 44 \\
11 & 1,024 & 63 \\
12 & 2,048 & 122 \\
13 & 4,096 & 190 \\
14 & 8,192 & 362 \\
15 & 16,384 & 612 \\
16 & 32,768 & 1,162 \\
\bottomrule
\end{tabular}
\endgroup
\end{table}

轨道数总和为 2,626。

\subsection{显式代表元集合相等}

有界低周期分类使用的不仅是式~\eqref{eq:burnside-orbit-count} 的计数。对每个 \(p\)，
任取前 \(p-1\) 项，再强制最后一项使乘积为正，从而构造全部 \(2^{p-1}\) 个
合法字。把每个字换成其所有旋转与反射后旋转中字典序最小者。将所得规范字
集合与存储表的 \((p,Q)\) 集合精确比较，并要求表标识符遵循该字典序。集合
相等同时排除遗漏轨道和用新标识符掩盖的重复轨道。

对每个存储代表元，第二项检查重新计算其二面体像集、轨道大小、周期提升、
本原 \(Q\) 周期和本原 \(\tau\) 周期。这证明周期 8 与 16 的两个目标行由周期胞
重复关联，且没有移除真正不同的本原相位。

\subsection{有限最小性覆盖}

直至 20 阶，直接迭代式~\eqref{eq:cycle-coordinate-space} 覆盖每个切换类。
在更大阶数上，按缺陷壳层生成规范 \((Q,\alpha)\) 记录。每条记录保存其二面
体轨道大小，从而保存式~\eqref{eq:orbit-weight-identity} 给出的代表重数。对每
个壳层，将这些重数之和与合法字的二项式奇偶计数比较；再把各壳层之和与
\(2^{n+1}\) 比较。末端游标证明规范生成器已经穷尽其有序定义域。

对每个 \(n=24,26,28,30\)，正式固定重量项链记录流都与独立 C 全空间扫描器
比较。后者遍历全部二进制字，检验合法奇偶性，直接构造二面体轨道，并在位图中
标记每个轨道成员；随后从磁盘映射表中消耗对应正式记录，独立比较规范字与轨道
大小。判定不依赖记录顺序或摘要相等。四阶分别恰好消耗 176,906、649,532、
2,405,236 和 8,964,800 条规范通量记录，没有缺失、重复、未消耗或轨道大小
不一致的记录。

每阶独立求得的轨道大小之和均为 \(2^{n-1}\)，缺陷直方图和轨道大小直方图也
逐项相同。添入两种和乐与两个全局取反提升后，恰好恢复 \(2^{n+1}\) 个切换类。
原有的 24 阶访问集合比较以及 Burnside、壳层、顺序和奇偶性核验仍作为附加
检查。由此在第~\ref{sec:computational-verification} 节披露的复核边界下建立
有限覆盖性。
